\section{Execução, Validação e Evidência}

Esta secção resume os pontos principais para validação do repositório. O procedimento completo de setup deve ser consultado em \repoPath{docs/local-baseline-setup.md}. Este documento não substitui esse guia.

\subsection{Requisitos gerais}

A execução local da baseline pode depender dos seguintes componentes:

\begin{itemize}
    \item .NET SDK compatível com a solução;
    \item Docker Desktop para os serviços locais;
    \item PowerShell para scripts de setup, runtime e evidência;
    \item Node.js e npm para a interface web;
    \item Git para clonagem e consulta da versão do repositório.
\end{itemize}

\subsection{Setup local}

A instrução principal para preparar a baseline local é:

\begin{center}
\repoPath{docs/local-baseline-setup.md}
\end{center}

Antes de executar comandos isolados, deve ser lido esse documento, porque ele concentra a sequência esperada de preparação do ambiente.

\subsection{Build e testes .NET}

Para validar a solução .NET, os comandos base são:

\begin{verbatim}
dotnet restore
dotnet build .\NatureProtector.sln --nologo -v minimal --configfile NuGet.Config
dotnet test .\NatureProtector.sln --no-restore --nologo -v minimal -m:1
\end{verbatim}

O build valida a compilação dos projetos. Os testes validam o comportamento do domínio, pipeline, infraestrutura, API, simulador e componentes partilhados.

\subsection{Interface web}

A interface web encontra-se em \repoPath{webUI/}. Para validar a compilação da UI:

\begin{verbatim}
cd webUI
npm install
npm run build
\end{verbatim}

Durante desenvolvimento, a UI é normalmente servida pelo Vite e comunica com a Backoffice API através do proxy configurado em \repoPath{webUI/vite.config.ts}. O arranque integrado da runtime deve ser feito pelos scripts de desenvolvimento indicados no setup.

\subsection{Infraestrutura local}

Os serviços locais são definidos por \repoPath{docker-compose.yml} e apoiados por scripts em \repoPath{infra/scripts/}. Em uso normal, deve ser seguido o procedimento descrito em \repoPath{docs/local-baseline-setup.md}. Quando aplicável, a infraestrutura pode ser iniciada por script próprio, por exemplo:

\begin{verbatim}
.\infra\scripts\up.ps1
\end{verbatim}

Os serviços locais podem incluir PostgreSQL, RabbitMQ, InfluxDB e Grafana.

\subsection{Runtime de desenvolvimento}

O arranque integrado da API, Prevention Host e UI pode ser feito por script de desenvolvimento, quando disponível:

\begin{verbatim}
.\scripts\dev\start-local-runtime.ps1 -OpenBrowser -ForceRestart
\end{verbatim}

Este script pressupõe que as dependências Docker estão disponíveis e que as portas usadas pela API e pela UI estão livres. O script gera evidência e sumários em \repoPath{docs/evidence/dev-runtime/}, quando configurado para isso.

\subsection{Execução de cenários e evidência}

Os cenários e smoke tests são apoiados por scripts em \repoPath{scripts/scenarios/} e \repoPath{scripts/evidence/}. Um exemplo de validação B/C pode ser executado em modo de simulação do próprio script:

\begin{verbatim}
.\scripts\evidence\run-v1-bc-smoke.ps1 -DryRun
\end{verbatim}

Quando a runtime local estiver ativa, a execução real pode apontar para a API:

\begin{verbatim}
.\scripts\evidence\run-v1-bc-smoke.ps1 -ApiBaseUrl http://localhost:5254
\end{verbatim}

A evidência gerada deve ser consultada em \repoPath{docs/evidence/}, em especial nas subpastas de runs, diagnostics e dev-runtime.

\subsection{Coverage}

A cobertura de testes pode ser gerada pelo script próprio, quando disponível:

\begin{verbatim}
.\scripts\tests\generate-coverage-report.ps1
\end{verbatim}

O output típico é guardado em \repoPath{coveragereport_core/}. Esta pasta é um artefacto gerado e não deve ser confundida com código fonte.
